\documentclass[11pt]{article}

\usepackage[margin=0.82in]{geometry}
\usepackage[T1]{fontenc}
\usepackage{lmodern}
\usepackage{microtype}
\usepackage{amsmath,amssymb,mathtools}
\usepackage{booktabs,tabularx,array}
\usepackage[ruled,vlined,linesnumbered]{algorithm2e}
\usepackage[dvipsnames]{xcolor}
\usepackage{enumitem}
\usepackage{xurl}
\usepackage{hyperref}

\hypersetup{
  colorlinks=true,
  linkcolor=MidnightBlue,
  citecolor=ForestGreen,
  urlcolor=BrickRed,
  pdftitle={When Should Personal Evidence Become Model Weights?}
}

\newcolumntype{Y}{>{\raggedright\arraybackslash}X}
\newcolumntype{L}[1]{>{\raggedright\arraybackslash}p{#1}}
\setlength{\tabcolsep}{5pt}
\renewcommand{\arraystretch}{1.13}
\setlist[itemize]{leftmargin=1.35em,itemsep=2pt,topsep=3pt}
\setlist[enumerate]{leftmargin=1.55em,itemsep=3pt,topsep=3pt}

\newcommand{\adapter}{\Delta\theta_u}
\newcommand{\memory}{M_u}
\newcommand{\baseparams}{\theta_0}
\newcommand{\prefscore}{U_{\mathrm{pref}}}
\newcommand{\taskscore}{U_{\mathrm{task}}}
\newcommand{\safescore}{U_{\mathrm{safe}}}

\title{\bfseries When Should Personal Evidence Become Model Weights?\\[3pt]
\large A Code-Grounded Method and Novelty Audit for\\
Per-User Offline-to-Online Personalization}
\author{Technical Design Note}
\date{August 29, 2026}

\begin{document}
\maketitle

\begin{abstract}
This note asks a narrow question: when should evidence about one user remain
in an editable prompt or memory, and when should it update that user's model
adapter?  The question matters because a parameter update is slower, harder to
inspect, and more likely to affect unrelated tasks.  It is useful only when it
improves future behavior beyond what the same evidence can achieve as context.

Two local codebases already contain real training.  One trains a Qwen3-4B
LoRA adapter with DPO.  The other implements GRPO and online policy
distillation (OPD).  These mechanisms are not new contributions.  Prior work
also covers per-user adapters, continual adapter updates, retrieval memory,
and context-conditioned steering.  The closest method, SPRInG, already
combines a user adapter with retrieved history.

The remaining hypothesis is a paired placement test.  For the same evidence,
build one context candidate and one adapter candidate.  Compare them on later
tasks, then activate the adapter only if it beats the context candidate while
preserving task success and safety.  This is a proposed experiment, not a
verified result or a settled novelty claim.
\end{abstract}

\section{Decision in Plain Language}

A prompt is the default place for new personal information.  It is easy to
read, change, delete, and limit to one task.  A user adapter is worth testing
only when a preference repeats across sessions, remains stable, applies to
several tasks, and should work without profile text at runtime.

Learning \(\adapter\) is not useful merely because the code can train it.
The update must answer a harder question:
\begin{quote}
Does the adapter improve later tasks more than giving the same evidence to the
model as prompt or memory, without reducing task success or safety?
\end{quote}
If the answer is no or remains uncertain, the evidence should stay in editable
context.

\section{Terms Used in This Note}

\begin{table}[ht]
  \centering
  \caption{Four objects that should not be treated as the same thing.}
  \label{tab:objects}
  \begin{tabularx}{\textwidth}{@{}L{0.17\textwidth}L{0.28\textwidth}YY@{}}
    \toprule
    Object & Meaning & Used when & Main property \\
    \midrule
    Profile prompt &
    A short text statement of user facts or preferences &
    Added directly to the input &
    Easy to inspect and edit, but consumes context \\
    Retrieved memory \(\memory\) &
    Selected records from the user's earlier interactions &
    Retrieved for a relevant request &
    Keeps details outside the model, but retrieval can miss or add noise \\
    Training hint \(h\) &
    A correction shown to a teacher during training &
    Used to make a target or distillation signal &
    A supervision channel, not necessarily deployed user state \\
    User adapter \(\adapter\) &
    A learned parameter change for user \(u\) &
    Loaded with the frozen base model &
    Works without profile text, but is harder to inspect and undo \\
    \bottomrule
  \end{tabularx}
\end{table}

\paragraph{Post-deployment online update.}
In this note, an online update occurs after the system has started serving the
user.  New user evidence arrives, an optimizer changes parameters, and the
system produces a new loadable checkpoint.  Training on a fixed dataset with
fresh model rollouts is often called online RL, but it is not the same as
learning from a deployed user's later interactions.

\paragraph{Personalized agent.}
A personalized agent completes tasks through tools, memory, actions, or
clarifying questions while adapting to one user.  A system that only changes
the style or content of generated text is personalized generation, not
necessarily personalized agent execution.

\section{What the Local Code Already Implements}

The two repositories contain training algorithms, not only data generation.

\begin{table}[ht]
  \centering
  \caption{Code-grounded status of the two local implementations.}
  \label{tab:code}
  \begin{tabularx}{\textwidth}{@{}L{0.18\textwidth}L{0.28\textwidth}YY@{}}
    \toprule
    Repository & Implemented update & Current evidence & Main defect \\
    \midrule
    \path{/fsx/xzhichao/workspace/offline training} &
    Qwen3-4B LoRA trained with \texttt{TRL DPOTrainer} &
    A clean comparison reports \(48.69\%\) macro and \(49.12\%\) micro win
    rate, with sign-test \(p=0.726\) &
    The profile is extracted before the holdout split, synthetic pairs are
    assumed to be preferences, and there is no adapter-only test arm \\
    \path{/fsx/xzhichao/workspace/oocso} &
    GRPO, hint-conditioned OPD, replay behavior cloning, and optimizer updates &
    The current three-user run is incomplete; an older result file uses a
    different configuration &
    The present offline phase uses scalar-reward GRPO rather than OPD, and no
    complete current run establishes an offline-to-online gain \\
    \bottomrule
  \end{tabularx}
\end{table}

The DPO result does not show a reliable training gain.  It also does not prove
that user adapters cannot work.  The comparison tested trained-plus-profile
against base-plus-profile, so it did not isolate what the adapter learned.

OOCSO already implements ``teacher sees a hint, student does not.''  A future
paper can use OPD as a training backend, but cannot present that mechanism as
new.

\section{Why Not Use a Prompt for Everything?}

\subsection{Cases where prompt or memory is better}

Keep evidence outside the weights when it is:
\begin{itemize}
  \item observed once or weakly supported;
  \item temporary, such as a tone request for one report;
  \item tied to one task, project, or tool;
  \item likely to change;
  \item easy to state as a short rule; or
  \item sensitive enough that the user should be able to inspect and delete it.
\end{itemize}

This choice has real costs.  Retrieval can select the wrong record.  Long
profiles consume context and are processed again for every request.  The model
may also ignore a relevant instruction.

\subsection{Cases where an adapter may help}

Test an adapter when the preference:
\begin{itemize}
  \item appears in several independent sessions;
  \item remains stable over time;
  \item affects several task families;
  \item cannot be expressed well as one short instruction; and
  \item should affect behavior even when profile text is absent.
\end{itemize}

An adapter may reduce repeated prompt cost and make the behavior more
consistent.  It also introduces training cost, forgetting risk, broader
behavior changes, and rollback work.  These costs make \(\adapter\) an
empirical choice rather than a default.

\section{Closest Prior Work}

\subsection{OPPU and Personalized DPO}

OPPU trains one PEFT module from each user's historical behavior
\cite{tan2024oppu}.  It can also use retrieved history or a generated profile.
It is an offline method.  It does not decide, for each later observation,
whether that observation should remain in memory or enter the adapter.

Personalized RLHF, evaluated through personalized DPO, learns user-conditioned
representations and LoRA parameters from offline preference pairs
\cite{li2024prlhf}.  It does not study chronological post-deployment updates,
retrieval memory, or agent task outcomes.

\subsection{SPRInG}

SPRInG is the closest personalization method \cite{kim2026spring}.  It already
contains:
\begin{itemize}
  \item one evolving LoRA adapter per user;
  \item likelihood-based selection of incoming interactions;
  \item a bounded buffer for examples that remain hard after an update;
  \item relevance-gated retrieval; and
  \item token-level interpolation between adapter-only and
  retrieval-conditioned generation.
\end{itemize}

Thus, ``combine memory and an adapter'' is not new.  SPRInG chooses training
examples through likelihood drift, then retains post-update residuals.  It
does not train matched context and adapter candidates from the same evidence
and select between them using later task success and safety.

SPRInG simulates sequential updates with chronological LongLaMP records.  It
personalizes text generation.  It does not execute tool-based agent tasks or
learn from live post-deployment agent outcomes.

\subsection{EvolveR}

EvolveR studies an agent-level experience loop rather than per-user
personalization \cite{wu2025evolver}.  It stores distilled principles as
retrieved context and updates a global policy with GRPO.  Its standard method
masks retrieved experience tokens from the training loss.

EvolveR also tests direct experience absorption by unmasking those tokens.
Its reported average drops from \(0.382\) to \(0.371\).  The authors attribute
the drop to noisy or irrelevant retrieved principles.  This result supports a
conservative design: context should not be written into weights without a
quality and relevance test.

\subsection{CoSteer}

CoSteer is a tuning-free alternative \cite{wu2025costeer}.  A local small model
runs with and without private context.  The difference between its two token
distributions steers a frozen cloud model.  CoSteer does not learn
\(\adapter\).  It shows that runtime personalization can be stronger than
plain prompting without changing model weights, so it is a useful baseline
when privacy and local inference matter.

\section{Novelty Audit}

\begin{table}[ht]
  \centering
  \caption{What the closest methods already cover.  ``Sequential'' means that
  later records can change model parameters.}
  \label{tab:coverage}
  \begin{tabularx}{\textwidth}{@{}L{0.16\textwidth}L{0.18\textwidth}L{0.19\textwidth}L{0.18\textwidth}Y@{}}
    \toprule
    Method & Learned object & External user state & Sequential update & Agent task outcomes \\
    \midrule
    OPPU &
    Per-user PEFT &
    Optional profile or retrieval &
    No &
    No \\
    Personalized DPO &
    Shared LoRA and user representation &
    Optional user text &
    No &
    No \\
    SPRInG &
    Per-user LoRA &
    Residual retrieval buffer &
    Yes, on a historical stream &
    No \\
    EvolveR &
    Global agent policy &
    Experience Base &
    Yes, during agent training &
    Yes, but not per user \\
    CoSteer &
    None &
    Private local context &
    No parameter update &
    No \\
    OOCSO / OPD &
    Per-user LoRA &
    Training hints &
    Implemented in simulation &
    Simulated rewards \\
    Proposed test &
    Per-user adapter when accepted &
    Editable memory otherwise &
    Yes, after deployment &
    Required \\
    \bottomrule
  \end{tabularx}
\end{table}

\begin{table}[ht]
  \centering
  \caption{The narrow mechanism-level difference under study.}
  \label{tab:delta}
  \begin{tabularx}{\textwidth}{@{}L{0.17\textwidth}L{0.25\textwidth}L{0.25\textwidth}Y@{}}
    \toprule
    Method & How evidence enters weights & How evidence remains external &
    How an update is accepted \\
    \midrule
    SPRInG &
    High likelihood-drift examples train the adapter &
    Hard post-update residuals enter the retrieval buffer &
    Fixed selection rule; no matched context-versus-adapter validation \\
    EvolveR &
    Agent trajectories update the policy; experience-token absorption is an
    ablation &
    Distilled principles are retrieved as context &
    Training reward; no per-evidence storage decision \\
    OOCSO / OPD &
    The no-hint student learns from a hint-conditioned teacher &
    Hints supervise training but are not a matched deployed memory arm &
    Reward-driven update; no context-versus-adapter activation test \\
    Proposed test &
    One candidate adapter is trained from eligible evidence &
    A matched candidate stores the same evidence as context &
    Later tasks must favor the adapter and pass task and safety constraints \\
    \bottomrule
  \end{tabularx}
\end{table}

\paragraph{What is not new.}
Per-user LoRA, personalized DPO, continual adapter updates, retrieval plus
parameters, hint distillation, and context-versus-parameter ablations all
exist.

\paragraph{What may remain distinct.}
The tentative difference is an evidence-level, matched comparison.  The same
evidence creates a context candidate and an adapter candidate from the same
parent checkpoint.  Later chronological tasks decide where the evidence
should live.  The adapter is accepted only when its advantage over context is
large enough and task and safety scores do not regress.

\paragraph{Current verdict.}
This is a defensible research hypothesis, not yet a publication novelty claim.
The claim becomes credible only after a broader search confirms that no prior
method performs this paired placement and constrained activation procedure.

\section{Proposed Paired Placement Test}

\subsection{State and evidence}

For user \(u\) at time \(t\), the active state is
\begin{equation}
  S_u^{(t)} =
  \left(\baseparams,\adapter^{(t)},\memory^{(t)},c_t\right),
\end{equation}
where \(c_t\) identifies the active checkpoint.  The frozen base parameters
are \(\baseparams\), the user adapter is \(\adapter^{(t)}\), and the editable
memory is \(\memory^{(t)}\).

Each evidence record is
\begin{equation}
  e_i=(u,t_i,x_i,\tau_i,y_i,f_i,o_i,p_i).
\end{equation}
Here \(x_i\) is the task, \(\tau_i\) is the agent trajectory, \(y_i\) is the
final answer, \(f_i\) is user feedback, \(o_i\) is the environment outcome,
and \(p_i\) records provenance.  Evidence cannot train a model when the user,
task, outcome, or generating checkpoint is unknown.

\subsection{Matched candidates}

Let \(E_t\) be a checked batch of new evidence.  Build both candidates from the
same active state.

\paragraph{Context candidate.}
\begin{equation}
  M_{u,\mathrm{ctx}} =
  \operatorname{MemoryUpdate}\left(\memory^{(t)},E_t\right),
  \qquad
  \Delta\theta_{u,\mathrm{ctx}}=\adapter^{(t)}.
\end{equation}

\paragraph{Adapter candidate.}
\begin{equation}
  \Delta\theta_{u,\mathrm{par}} =
  \operatorname{TrainAdapter}\left(\adapter^{(t)},E_t\right),
  \qquad
  M_{u,\mathrm{par}}=\memory^{(t)}.
\end{equation}
The training backend may be SFT, DPO, or OPD.  These are inherited
components.  The adapter candidate is evaluated without adding the new
evidence to its prompt.

\subsection{Activation rule}

Let \(V_t\) contain sessions that occur after \(E_t\) but before the untouched
final test period.  Evaluate the context candidate \(m_{\mathrm{ctx}}\), the
adapter candidate \(m_{\mathrm{par}}\), and the current model \(m_0\) under the
same decoding settings.

Define the paired personalization difference
\begin{equation}
  D_{\mathrm{pref}} =
  \prefscore(m_{\mathrm{par}};V_t)
  -
  \prefscore(m_{\mathrm{ctx}};V_t).
\end{equation}
Activate the adapter only if
\begin{align}
  \operatorname{LCB}(D_{\mathrm{pref}}) &> \delta, \label{eq:prefgate}\\
  \operatorname{LCB}\!\left(
    \taskscore(m_{\mathrm{par}};V_t)
    -
    \taskscore(m_0;V_t)
  \right) &\ge -\epsilon, \label{eq:taskgate}\\
  \safescore(m_{\mathrm{par}};V_t)
    &\ge \safescore(m_0;V_t). \label{eq:safetygate}
\end{align}
The lower confidence bound (LCB) prevents a small noisy difference from
triggering an update.  Equation~\eqref{eq:taskgate} is a non-inferiority test:
task success may fall by at most the predeclared margin \(\epsilon\).

If the adapter fails, use context only when context improves personalization
without reducing task or safety scores.  Otherwise, keep the evidence pending,
ask the user for clarification, or make no change.  The procedure never forces
every record into memory or weights.

\begin{algorithm}[ht]
\caption{Paired evidence placement and conservative activation}
\label{alg:placement}
\KwIn{Active state \(S_u^{(t)}\), new evidence \(E_t\), later validation
window \(V_t\), thresholds \(\delta,\epsilon\)}
\KwOut{Versioned next state \(S_u^{(t+1)}\)}
Check user identity, time, task outcome, feedback, and checkpoint provenance\;
Remove unsupported, contradictory, unsafe, or task-failing records\;
\If{eligible evidence is insufficient}{
  keep the active state and record a pending decision\;
  \Return \(S_u^{(t)}\)\;
}
Build \(m_{\mathrm{ctx}}\) by storing \(E_t\) as editable context\;
Build \(m_{\mathrm{par}}\) by training a child adapter on the same \(E_t\)\;
Evaluate \(m_0\), \(m_{\mathrm{ctx}}\), and \(m_{\mathrm{par}}\) on the same \(V_t\)\;
\eIf{Eqs.~\eqref{eq:prefgate} to \eqref{eq:safetygate} pass}{
  activate the child adapter and save \(m_0\) as the rollback target\;
}{
  \eIf{the context candidate improves personalization and passes task and
  safety checks}{
    activate the context candidate\;
  }{
    make no change or ask the user to clarify\;
  }
}
Write the data boundary, metrics, parent, child, and decision to the ledger\;
\Return the versioned next state\;
\end{algorithm}

\section{Experiment That Can Answer the Question}

\subsection{Chronological data boundary}

Split whole sessions before extracting a profile or generating training data:
\begin{equation}
  H_u^{\mathrm{past}}
  \prec E_t
  \prec V_t
  \prec H_u^{\mathrm{test}}.
\end{equation}
The order symbol means ``occurs earlier than.''  The profile, memory, hints,
preference pairs, and adapter may use only the records available at that time.
The final test period remains untouched until the method and thresholds are
fixed.

\subsection{Required arms}

\begin{table}[ht]
  \centering
  \caption{Minimum comparison needed to identify where personalization comes
  from.}
  \label{tab:arms}
  \begin{tabularx}{0.94\textwidth}{@{}L{0.31\textwidth}Y@{}}
    \toprule
    Arm & Question answered \\
    \midrule
    Frozen base &
    What happens without personalization? \\
    Base plus profile or memory &
    How much can explicit context do without training? \\
    User adapter only &
    Did the weights store useful personal behavior? \\
    User adapter plus memory &
    Are the two representations complementary? \\
    OPPU-style adapter &
    Does direct offline per-user training suffice? \\
    SPRInG &
    Does likelihood-based selective adaptation plus retrieval suffice? \\
    OOCSO-style OPD &
    Does hint distillation improve later no-hint behavior? \\
    CoSteer, when deployment permits &
    Can tuning-free local steering avoid a parameter update? \\
    Paired placement test &
    Does evidence-level placement improve over fixed storage rules? \\
    \bottomrule
  \end{tabularx}
\end{table}

\subsection{Separate outcomes}

Do not combine all outcomes into one score.  Report:
\begin{itemize}
  \item preference match on future sessions;
  \item environment-verified task success;
  \item safety violations and invalid actions;
  \item unrelated-task regression;
  \item adaptation speed after a real preference change;
  \item forgetting after later updates;
  \item prompt tokens, training cost, and inference cost;
  \item update count, rejection count, and rollback count; and
  \item uncertainty intervals over users and tasks.
\end{itemize}

Pairwise judges must allow ties and invalid outputs.  A style judge cannot
replace task correctness.  A result from one synthetic user cannot establish a
general personalization method.

\section{Implementation Order}

The shortest credible path is:
\begin{enumerate}
  \item Split the offline DPO conversations by time before profile extraction.
  \item Add base, base-plus-profile, adapter-only, and
  adapter-plus-profile evaluation arms.
  \item Verify every synthetic pair with separate preference, task, and safety
  checks.
  \item Reproduce one clean offline result.
  \item Finish one current OOCSO run with fixed configuration and complete
  checkpoint lineage.
  \item Add the evidence ledger and the two matched candidates around the
  existing SFT, DPO, or OPD backend.
  \item Test placement on controlled stable, temporary, task-local,
  contradictory, and drifting preferences.
  \item Move to real tool-using agent trajectories only after the storage
  decision works on the smaller test.
\end{enumerate}

\section{Final Assessment}

The local projects already contain DPO, GRPO, OPD, LoRA updates, and
evaluation drivers.  They are not data-generation-only projects.  Their
current results do not show that a user adapter reliably improves future
behavior.

There is also no general reason to prefer \(\adapter\) over a prompt.  Prompt
or memory is better for recent, uncertain, temporary, task-local, changing, or
sensitive evidence.  An adapter is worth the cost only for repeated and stable
behavior that transfers across tasks and works without runtime profile text.

The proposed paired placement test is narrower than ``memory plus adapter.''
It asks both representations to compete using the same evidence and later
tasks.  SPRInG, EvolveR, and the local OPD implementation make this a plausible
next experiment, but they also narrow the novelty claim.  The paper should
claim a new method only after the literature audit and matched experiment both
support that claim.

\begin{thebibliography}{9}

\bibitem{tan2024oppu}
Zhaoxuan Tan, Qingkai Zeng, Yijun Tian, Zheyuan Liu, Bing Yin, and Meng Jiang.
\newblock Democratizing Large Language Models via Personalized
Parameter-Efficient Fine-Tuning.
\newblock In \emph{EMNLP}, 2024.
\newblock \url{https://aclanthology.org/2024.emnlp-main.372/}.

\bibitem{li2024prlhf}
Xinyu Li, Ruiyang Zhou, Zachary C. Lipton, and Leqi Liu.
\newblock Personalized Language Modeling from Personalized Human Feedback.
\newblock arXiv:2402.05133, 2024.
\newblock \url{https://arxiv.org/abs/2402.05133}.

\bibitem{kim2026spring}
Seoyeon Kim and Jaehyung Kim.
\newblock SPRInG: Continual LLM Personalization via Selective Parametric
Adaptation and Retrieval-Interpolated Generation.
\newblock arXiv:2601.09974, 2026.
\newblock \url{https://arxiv.org/abs/2601.09974}.

\bibitem{wu2025evolver}
Rong Wu, Xiaoman Wang, Jianbiao Mei, and others.
\newblock EvolveR: Self-Evolving LLM Agents through an Experience-Driven
Lifecycle.
\newblock arXiv:2510.16079, 2025.
\newblock \url{https://arxiv.org/abs/2510.16079}.

\bibitem{wu2025costeer}
Hang Lv, Sheng Liang, Hao Wang, and others.
\newblock CoSteer: Collaborative Decoding-Time Personalization via Local
Delta Steering.
\newblock arXiv:2507.04756, 2025.
\newblock \url{https://arxiv.org/abs/2507.04756}.

\end{thebibliography}

\end{document}
